\documentclass{article}

% No \usepackage commands – everything here is part of plain LaTeX

\title{A Sample Document for Testing GitHub Pages}
\author{Test User}
\date{\today}

\begin{document}

\maketitle

\begin{abstract}
This document is meant to test the LaTeX compilation pipeline on GitHub Pages.
It includes sections, mathematical expressions, a table, a figure, cross‑references,
citations, and a bibliography — all using only built‑in LaTeX features.
\end{abstract}

\section{Introduction}
This is a simple test document. It shows how to structure a paper without relying on
external packages. We can cite a well‑known book~\cite{knuth:1984} or a research article~\cite{lamport:1994}.

\section{Mathematical Content}
Mathematics can be written inline, like $E = mc^2$, or displayed:

\[
\int_{-\infty}^{\infty} e^{-x^2} dx = \sqrt{\pi}.
\]

We can also number equations using the \texttt{equation} environment:
\begin{equation}
\label{eq:identity}
\sum_{i=1}^{n} i = \frac{n(n+1)}{2}.
\end{equation}
Equation~\eqref{eq:identity} is a classic identity.

\section{Tables and Lists}
A simple table:
\begin{center}
\begin{tabular}{|l|c|r|}
\hline
\textbf{Item} & \textbf{Quantity} & \textbf{Price} \\
\hline
Apple & 3 & \$1.50 \\
Orange & 2 & \$1.00 \\
\hline
\end{tabular}
\captionof{table}{Fruit prices}  % \captionof is not standard; use \caption inside table environment
\end{center}

Instead, we place the table inside a \texttt{table} environment to get a caption:
\begin{table}[h]
\centering
\begin{tabular}{|l|c|r|}
\hline
\textbf{Item} & \textbf{Quantity} & \textbf{Price} \\
\hline
Apple & 3 & \$1.50 \\
Orange & 2 & \$1.00 \\
\hline
\end{tabular}
\caption{Fruit prices}
\label{tab:prices}
\end{table}

Table~\ref{tab:prices} shows some sample data.

An itemised list:
\begin{itemize}
\item First point.
\item Second point, with a sublist:
  \begin{enumerate}
  \item Subitem one.
  \item Subitem two.
  \end{enumerate}
\item Third point.
\end{itemize}

\section{Figures}
We can include a figure without any external graphics package by using the \texttt{picture} environment or simply a \texttt{\string\rule} to draw a box:
\begin{figure}[h]
\centering
\rule{4cm}{2cm}  % a black rectangle
\caption{A simple figure made with \texttt{\string\rule}.}
\label{fig:rule}
\end{figure}
Figure~\ref{fig:rule} shows a placeholder graphic.

\section{Conclusion}
This document has demonstrated several basic LaTeX features. For more advanced
typography, one would normally load packages like \texttt{amsmath}, \texttt{graphicx}, etc.
However, even with only the built‑in commands, a well‑structured document can be produced.

\bibliographystyle{plain}
\bibliography{references}

\end{document}